\begin{filesystemlayout}

\chapter{文件系统的形成：从可寻址空间到持久对象}

\begin{keyidea}
文件系统并不是从 inode、超级块或目录树开始的。最初只有一片能够在重启后保留内容的
可寻址空间；只有在多个对象共享这片空间之后，系统才依次需要管理单位、分配状态、
对象身份、名称关系和重新挂载的入口。本章沿着这条需求链建立最小磁盘格式，并说明
每一种结构为何出现、保存什么事实，以及它不能替代什么。
\end{keyidea}

\section{从一个块开始：让问题逐个逼出文件系统概念}

这一节故意不从 inode、超级块或 VFS 开始。直接给出这些名词，读者只能背定义，
却不知道它们为什么存在。我们只先接受一个事实：掉电后还想保留的数据，必须
最终进入某种非易失存储设备；操作系统把这类设备先抽象成一串可编号的块。

这里的\emph{块（block）}是一次寻址和管理所使用的固定大小区间。假设块大小为
4096 字节，那么块 7 表示从设备字节偏移 $7\times4096$ 开始的 4096 字节。
选择固定大小不是为了让文件也固定大小，而是为了让“第几块”可以用一个整数表示，
让分配、缓存、校验和 I/O 请求都有明确边界。设备逻辑块、文件系统块和内存页
可能大小不同；本节统一取 4096 字节，只是为了让推导可以手算。

\begin{keyidea}
块只回答两个问题：\textbf{位置在哪里、一次管理多少字节}。它还没有回答
“这些字节属于谁”“有效长度是多少”“叫什么名字”“重启后从哪里开始找”。
文件系统中的其他概念，都是为了补上这些缺失的信息。
\end{keyidea}

\subsection{第一步：一个块为什么还不能叫文件}

假设设备现在只有一个 4096 字节块，我们把字符串 \code{"hello"} 写在块首部。
重启后读出这个块，会立刻遇到三个问题：

\begin{enumerate}\tightlist
\item 有效内容究竟是 5 字节，还是后面的零也属于内容？
\item 这 5 字节是文本、目录、图片，还是某个程序的内部状态？
\item 如果写到一半掉电，怎样判断块里的内容能不能信任？
\end{enumerate}

因此，只有数据字节还不够。系统还需要一些\emph{描述数据的数据}，例如类型、
有效长度和校验信息。描述其他数据的数据叫作\emph{元数据（metadata）}。
“元”不是神秘的新介质；它仍然是普通字节，只是解释角色不同。

最小记录可以先写成：

\begin{lstlisting}
+----------+----------+----------------------+------------------+
| type     | length   | payload              | unused           |
| 1 byte   | 4 bytes  | length bytes         | rest of block    |
+----------+----------+----------------------+------------------+
\end{lstlisting}

有了 \code{length=5}，重启后才知道只应返回 \code{"hello"}。这一步得到的
不是完整文件系统，只是一个“带边界的持久记录”。

\subsection{第二步：第二个文件逼出身份和分配问题}

现在要同时保存 \code{a.txt} 和 \code{b.txt}。若程序各自随便挑一个块，
两个程序可能都选择块 20，后写者会覆盖先写者。因此必须有一个全局规则回答：

\begin{quote}
哪些块已经属于某个对象，哪些块仍可分配？
\end{quote}

最直接的办法是为每个块保存一位。位只有 0 和 1：0 表示空闲，1 表示已占用。
把这些位按块号排列，就得到\emph{块位图（block bitmap）}。例如：

\begin{lstlisting}
block number:  0 1 2 3 4 5 6 7 8 9 ...
bitmap bit:    1 1 1 1 0 0 1 0 0 0 ...

解释：块 0--3 和块 6 已被占用；下一次分配可以从值为 0 的位置中选择。
\end{lstlisting}

位图解决的是\emph{空间所有权冲突}，不是文件内容定位。看到“块 6 已用”，
仍不知道它属于哪个文件，也不知道文件有多长。

接着考虑一个 9000 字节文件。一个 4096 字节块放不下它，至少需要三个块；
这三个块还可能是 20、21、90，并不连续。系统于是需要为每个文件保存一份记录：

\begin{itemize}\tightlist
\item 文件的类型和权限；
\item 文件的精确字节长度；
\item 文件内容所在的块号，或能找到这些块的索引；
\item 链接计数、时间戳等生命周期信息。
\end{itemize}

这份“每个文件一份、描述文件但不保存文件名”的持久记录，就是
\emph{inode（index node，索引节点）}。这个名字强调它首先是一个索引入口：
给定 inode，就能找到该文件的数据块和元数据。

\begin{warningbox}
inode 不是文件内容，也不是文件名，更不是进程里的 fd。inode 是文件系统内部
对一个持久对象的身份记录。文件数据可以搬到别的块，名字也可以改变，但只要
对象仍是同一个对象，它的 inode 身份就可以保持不变。
\end{warningbox}

多个 inode 也需要分配和回收，所以还需要\emph{inode 位图}。块位图回答
“哪些数据块空闲”，inode 位图回答“哪些 inode 编号空闲”。两张位图管理的是
不同资源，不能合并成一张含义模糊的表。

\subsection{第三步：文件名为什么不直接塞进 inode}

用户希望通过 \code{/notes/today.txt} 找文件，而不是记住 inode 37。
因此必须有一种持久映射，把一个名字转换成 inode 编号：

\begin{lstlisting}
"today.txt"  -> inode 37
"todo.txt"   -> inode 52
\end{lstlisting}

保存这种映射的对象叫\emph{目录（directory）}；映射中的一项叫
\emph{目录项（directory entry，简称 dirent）}。目录不是磁盘之外的特殊容器，
它本身也是一种文件：普通文件的数据是用户字节，目录文件的数据是一组
“名字 $\rightarrow$ inode 编号”记录。

文件名不直接放在 inode 中有三个直接收益：

\begin{enumerate}\tightlist
\item 同一个 inode 可以被多个目录项引用，这就是硬链接；
\item 改名主要修改目录项，不必搬动文件内容；
\item inode 表可以保持固定大小记录，目录项则按名字长度独立组织。
\end{enumerate}

这里也解释了为什么删除名字不一定立刻删除数据。\code{unlink} 删除的是一个
目录项，并把 inode 的链接计数减一；只有链接计数已经为零，并且也没有进程
继续打开该文件时，inode 和数据块才真正具备回收条件。

\subsection{第四步：目录树从哪里开始，逼出根目录}

目录项能够从名字找到 inode，但路径解析必须先有起点。解析
\code{/notes/today.txt} 时，系统需要先找到 \code{/}，再在其中查
\code{notes}，最后在 \code{notes} 目录中查 \code{today.txt}。

这个固定起点叫\emph{根目录（root directory）}。根目录仍然有自己的 inode；
“根”只表示路径解析从它开始，并不表示它采用另一种磁盘格式。因此文件系统必须
在某个全局位置记住根目录 inode 编号。

\subsection{第五步：重启后的第一个块怎么找，逼出超级块}

现在磁盘上已经有 inode 位图、块位图、inode 表、目录数据和普通文件数据。
但机器重启后，内存中的所有指针都消失了。挂载程序面对的仍只是一串块：

\begin{quote}
块有多大？共有多少块？inode 表从哪一块开始？根目录 inode 是多少？
这个设备上的字节甚至是不是我们认识的文件系统？
\end{quote}

这些问题不能再通过 inode 回答，因为此时连 inode 表在哪里都不知道。必须规定
一个预先约定的位置，放置足以解释整个文件系统布局的全局元数据。这个入口记录
叫\emph{超级块（superblock）}。

“超级”表示它描述的是整个文件系统，而 inode 只描述单个文件。超级块至少需要：

\begin{itemize}\tightlist
\item magic：用于拒绝把随机字节误认成这种文件系统；
\item version：说明磁盘格式版本，避免用错误规则解释旧数据；
\item block size 和 block count：确定块的粒度和合法范围；
\item inode count：确定合法 inode 编号范围；
\item inode/块位图以及 inode 表的位置；
\item data start：数据区从哪里开始；
\item root inode：路径解析的固定起点。
\end{itemize}

magic 不是安全校验，也不能证明数据没有损坏；它只是格式身份标记。version 也不是
程序版本号，而是\emph{磁盘格式}版本。真正的损坏检测还需要字段范围检查、校验和、
副本或日志，这些机制后面再引入。

\begin{pdfdiagram}
  \node[os title] at (0,4.6) {挂载时的引导链：每一步只使用上一步刚获得的信息};
  \node[os control, minimum width=8.8cm] (dev) at (0,3.7) {已知事实：设备可按块号读取，超级块位于约定位置};
  \node[os memory, minimum width=8.8cm] (sb) at (0,2.6) {读超级块：验证 magic/version/范围，得到各区域位置};
  \node[os compute, minimum width=8.8cm] (itable) at (0,1.5) {按 root inode 计算 inode 表中的字节位置，读根 inode};
  \node[os state, minimum width=8.8cm] (root) at (0,0.4) {按根 inode 的块映射读取根目录项};
  \node[os small block, minimum width=8.8cm] (path) at (0,-0.7) {从根目录开始逐段解析路径，最终找到目标 inode};
  \draw[os strong bus] (dev) -- (sb);
  \draw[os strong bus] (sb) -- (itable);
  \draw[os strong bus] (itable) -- (root);
  \draw[os strong bus] (root) -- (path);
\end{pdfdiagram}
\diagramnote{这条链证明超级块的必要性：若没有一个位置固定、格式可验证的全局入口，重启后的系统就无法知道 inode 表和根目录在哪里，后续所有名字和文件都无从查找。}

\subsection{六个概念各自只解决一个核心问题}

\begin{longtable}{L{0.16\linewidth}L{0.30\linewidth}L{0.22\linewidth}L{0.24\linewidth}}
\toprule
\textbf{概念} & \textbf{没有它时首先坏在哪里} & \textbf{它保存什么} & \textbf{它不负责什么} \\
\midrule
块 & 无法用整数表达稳定位置和管理粒度 & 固定大小的字节区间 & 不表达文件身份和名字 \\
块位图 & 两个对象可能分到同一块 & 每个块是否已分配 & 不说明块属于哪个文件 \\
inode & 文件变长或不连续后，无处保存块映射 & 单个文件的身份、长度、类型和块映射 & 通常不保存文件名 \\
inode 位图 & 无法判断哪个 inode 编号可以复用 & inode 编号的分配状态 & 不保存 inode 正文 \\
目录项 & 用户只能记 inode 编号，无法使用路径 & 名字到 inode 编号的映射 & 不保存普通文件内容 \\
超级块 & 重启后不知道怎样解释磁盘布局 & 全局格式、区域位置和根 inode & 不保存每个文件的全部元数据 \\
\bottomrule
\end{longtable}

这张表也是后续阅读的约束：每当出现一个新字段或新对象，都必须指出它在解决哪一个
此前无法解决的问题。如果只是给出名词而没有失败场景，就还没有完成解释。

\subsection{一个最小磁盘布局怎样闭环}

教学文件系统 MiniFS-Lab 使用 64 个块，每块 4096 字节，总容量 256 KiB。
它只为了暴露机制，不追求容量和性能：

\begin{lstlisting}
block 0      superblock：全局格式与区域位置
block 1      inode bitmap：哪些 inode 编号已分配
block 2      block bitmap：哪些块已分配
block 3--6   inode table：32 个固定大小 inode
block 7      root directory data：根目录的目录项
block 8--62  ordinary data blocks
block 63     reserved：后续日志实验使用
\end{lstlisting}

为什么把超级块放在块 0？不是因为所有真实文件系统都必须如此，而是教学格式做了
这个约定，使 mount 知道第一步读哪里。真实文件系统可能预留引导区、在固定偏移放
超级块并保存备份；关键不在具体块号，而在“入口位置必须由格式事先约定”。

为什么 inode 表使用固定大小记录？因为给定 inode 编号 $i$ 后，可以直接计算：

\[
\text{byte offset}=\text{inode table start}+i\times\text{inode size}.
\]

这让读取 inode 不必先扫描所有旧 inode。代价是每个 inode 能内联保存的字段有限；
更大的可变信息需要放到其他块中。

为什么根目录先占用块 7？因为格式化完成后必须立刻存在一个路径起点。mkfs 会同时：

\begin{enumerate}\tightlist
\item 写超级块；
\item 在 inode 位图中标记根 inode 已用；
\item 在块位图中标记元数据块和块 7 已用；
\item 初始化根 inode，使它指向块 7；
\item 在块 7 写入 \code{"."} 和 \code{".."} 目录项。
\end{enumerate}

如果只写超级块而不初始化根 inode，文件系统虽然“能识别”，却没有路径起点；
如果初始化根 inode但忘记标记块位图，后续分配器可能把根目录块再次分给普通文件。
这说明 mkfs 不是打印一个 magic，而是在磁盘上建立一组互相一致的不变量。

\subsection{第一次创建文件：每个对象为什么都必须参与}

现在执行 \code{create("/hello", "abc")}。暂不考虑崩溃，最小正确顺序如下：

\begin{enumerate}\tightlist
\item 从超级块得到 root inode 编号和各区域位置；否则不知道去哪里解析路径。
\item 读取根 inode，再读取根目录块，确认名字 \code{hello} 尚不存在；否则会无意覆盖已有名字。
\item 扫描 inode 位图，选择一个空闲 inode 并标记已用；否则新文件没有稳定身份。
\item 扫描块位图，选择一个空闲数据块并标记已用；否则文件没有独占存储空间。
\item 初始化新 inode：类型为普通文件，长度为 3，数据块指向刚分配的块；否则读路径无法解释数据。
\item 把 \code{abc} 写入数据块；这是用户真正要求保存的内容。
\item 在根目录中加入 \code{"hello" -> new inode}；否则文件存在但没有路径可达。
\end{enumerate}

这里故意把“为什么”写进每一步。创建文件不是“写一份 inode”这么简单，而是同时修改
两张分配表、一个 inode、一个目录和一个数据块。也正因为一次逻辑操作对应多次块写，
掉电可能留下半完成状态；日志和 copy-on-write 将在后文由这个具体失败自然引出。

\begin{warningbox}
在学习文件系统时，必须始终区分三类关系：\textbf{位图说明资源是否被占用，
inode 说明文件数据在哪里，目录项说明什么名字能找到该 inode}。三者任何两个一致、
第三个不一致，都可能造成块重复分配、空间泄漏、悬空名字或不可达 inode。
\end{warningbox}

\subsection{磁盘记录与内存对象不是同一个东西}

磁盘 inode 是固定格式的持久字节；内核中的 \code{struct inode} 是运行时对象，
除了从磁盘读出的字段，还包含锁、引用计数、page cache 入口和操作表。磁盘目录项是
持久的名字映射；内存 dentry 是路径查找缓存。超级块也同时有磁盘格式和内存实例。

这样区分的原因是：磁盘格式必须跨重启稳定，不能随编译器结构体布局变化；内存对象
只服务当前运行，可以包含指针、锁和缓存。把 C++ 结构体直接 \code{memcpy} 到磁盘，
会受到字节序、对齐、填充和版本变化影响。MiniFS-Lab 因此使用显式的定宽整数编码，
并在 mount 时逐字段验证范围。

至此，超级块、inode、目录项和位图都不是凭空给出的名词，而是一条完整推导链：

\begin{lstlisting}
固定位置需求 -> block
多个对象不能冲突 -> allocation bitmap
变长、非连续文件需要身份和块映射 -> inode
人需要用名字寻找对象 -> directory entry
路径解析需要固定起点 -> root inode
重启后需要知道整个布局 -> superblock
多块更新可能只完成一半 -> journal / copy-on-write（后文）
\end{lstlisting}

后续各节会对这些概念做第二遍深化：先看真实 I/O 怎样把块送到设备，再看 ext4
inode、extent、目录索引、page cache 和日志怎样把同一组问题扩展到现代系统。

\subsection{可运行证据：MiniFS-Lab 不是纸上结构图}

完整参考实现位于
\filepath{source/latex/listings/minifs\_lab.cpp}，并由本书 \code{make check}
自动编译运行。它没有把 C++ 结构体直接复制到磁盘，而是逐字段写入小端定宽整数；
这样读者可以把每一个磁盘字节与本节格式表对应起来。

实现包含三组验收：

\begin{enumerate}\tightlist
\item \textbf{格式化与重挂载}：mkfs 写入超级块、两张位图、根 inode 和根目录项；
  mount 丢弃所有旧内存对象，只凭磁盘字节重新找到 \code{/hello}。
\item \textbf{删除与资源回收}：unlink 同时移除目录项、清 inode 位和块位，随后
  audit 验证“可达对象、位图、inode 块映射”三者一致。
\item \textbf{掉电注入}：create 已经分配 inode 和数据块、也写好了内容，但在发布
  目录项之前模拟掉电。重挂载后名字不存在，audit 却会发现一个不可达 inode 和一个
  无可达所有者的已分配块；修复器扫描可达关系后才能回收它们。
\end{enumerate}

\begin{lstlisting}
$ make -C books/operating-systems-volume-1 check
CXX source/latex/listings/minifs_lab.cpp
MiniFS-Lab checks passed
\end{lstlisting}

这个失败实验正是后文引入日志的理由。若 create 的多次块写没有事务边界，mount
只能看到“磁盘现在是什么样”，无法知道半完成操作原本想做什么；fsck 可以扫描并
修补结构，却不一定能恢复用户原本想发布的名字和内容。日志增加的不是另一个名词，
而是可判定的提交证据：恢复程序据此区分“应当重放的完整事务”和“应当丢弃的半成品”。


%% 文件系统：从磁盘字节到现代设计
%% Sections 1--4

%%% 下面这段设备级材料保留在源稿中供后续实践卷整理，不进入本册正文。
%%% 块层、DMA 与设备完成语义已经由前两章负责，本章从文件系统管理单位继续。
\iffalse

\fsdivision{块设备提供什么持久化能力}

\fstopic{教学块号与设备 I/O 的边界}

上一节从“一个块”推导出了位图、inode、目录项和超级块。那一遍有意把设备
简化成可读写的块数组；现在把镜头下移，回答这些块怎样经过控制器、DMA 和
Linux 块层成为真实 I/O。概念推导解释“为什么需要这些对象”，本节则解释
“对象中的块号最终怎样变成设备命令”。

从操作系统块层看，它可以先抽象成\emph{一个有固定容量、可按逻辑块寻址的线性空间}。这里的“块”不是宇宙通用的 4096 字节：设备逻辑块、设备物理块、文件系统块和内存页是四个相关但不同的粒度。为便于后面的算例，本节暂时假设文件系统块为 4096 字节。

\begin{lstlisting}
FS Block 0  FS Block 1  FS Block 2  ...     FS Block N-1
+-----------+-----------+-----------+-------+-----------+
|  4096 B   |  4096 B   |  4096 B   |  ...  |  4096 B   |
+-----------+-----------+-----------+-------+-----------+
     ^
     |
  块号 = 起始字节地址 / 4096
\end{lstlisting}

图中的 4096 字节只是本章示例。实际设备会报告 logical block size、physical block size、minimum I/O size 和 optimal I/O size；文件系统再选择自己的块大小，内存管理也有独立的页大小。CPU 不能直接用
\code{mov} 指令读写这些块——磁盘不在内存地址空间中。CPU 必须通过
\emph{磁盘控制器}（disk controller）间接访问。

\fstopic{DMA 路径：CPU 如何让磁盘干活}

现代高性能块设备通常用 \emph{直接内存访问}（Direct Memory Access,
DMA）传输数据。下面是队列式设备的简化路径；具体设备可能使用 submission queue、descriptor ring 或厂商自定义命令格式，不能把某一种寄存器布局当成所有设备的标准。

\begin{enumerate}\tightlist
\item 内核分配并 pin 住传输缓冲区，通过 DMA API 把它映射成设备可见的 DMA 地址；启用 IOMMU 时，这通常是 IOVA，不等同于 CPU 物理地址。
\item 驱动在内存队列或描述符环中准备命令和散布-收集列表，并在发布前完成必要的内存排序。
\item 驱动更新队列尾指针或写 doorbell MMIO 寄存器，通知控制器有新命令。
\item 控制器开始工作：它自行读取 scatter-gather 描述符，然后通过总线
  在设备与已映射缓冲区之间传输数据。传输本身不依赖 CPU 执行逐字节复制指令，但 CPU 仍可能在提交前填充写缓冲区、在完成后解析或复制数据。
\item 数据传输完成后，控制器向 CPU 发出一个中断（现代设备用 MSI-X
  消息信号中断）。
\item CPU 的中断处理程序被调用，确认 I/O 完成。
\end{enumerate}

在 DMA 传输期间，发起 I/O 的线程通常不必占用 CPU 等待；CPU 可以执行其他线程，
完成队列则可按设备与负载采用中断或轮询。DMA 规定数据怎样搬运，中断与轮询规定
完成怎样通知，两者处于不同层次。

\begin{lstlisting}
  CPU               Memory                Disk Controller        Disk
   |                  |                        |                  |
   |--setup SG list-->| (write to RAM)         |                  |
   |--publish queue entry---------------------->|                  |
   |--ring doorbell---------------------------->|                  |
   |                  |                        |--read SG list--->| (DMA)
   |     (free to     |<======== DMA transfer =====================>| (data)
   |      run other   |                        |                  |
   |      threads)    |                        |                  |
   |<-----------------|-------------------------|--MSI-X interrupt |
   |                  |                        |                  |
\end{lstlisting}

\fstopic{一个能工作的裸磁盘程序}

假设我们有一个简单的单任务环境。一个程序想把 8192 字节的数据存到磁盘上，
过一会儿再读回来。它可以这样设计：

\begin{lstlisting}[language=C]
// 写数据到磁盘：块 100--101，共 8192 字节
int fd = open("/dev/sda", O_RDWR);
lseek(fd, 100 * 4096, SEEK_SET);   // 定位到块 100
write(fd, my_data, 8192);            // 写 2 个块

// ... 时间流逝 ...

// 读回来
lseek(fd, 100 * 4096, SEEK_SET);
read(fd, my_data, 8192);
\end{lstlisting}

这在单任务、单程序、文件数量很少的场景下\emph{可以工作}。程序自己记住
"数据在块 100--101"，就像你在仓库里记住"东西放在第 100 号货架上"。

但只要场景稍微复杂一点，问题就来了。

\fstopic{第一遍推导的设备级复核}

下表列出了裸磁盘模型在面对真实多任务、持久化需求时的核心问题：

\begin{longtable}{|F{0.14\linewidth}|F{0.58\linewidth}|F{0.22\linewidth}|}
\hline
\rowcolor{BookBlueTint}
\textbf{问题} & \textbf{具体困难} & \textbf{举例} \\\tablerowrule
\endhead
命名 & 块只有编号，没有名字。程序必须自己记住"块 100 是我的数据"。用户不可能记住每个文件在哪个块。 & \code{open("/home/alice/report.pdf")} \\\tablerowrule
分配 & 谁来跟踪哪些块是空闲的？程序 A 用了块 100，程序 B 不能再用。没有全局的空闲块记录。 & 两个程序同时选了块 100 \\\tablerowrule
增长 & 文件写完后可能还要追加数据。如果块 100 后面的块 101 已被别人占用，文件无法连续增长。 & \code{append()} 导致需要新块 \\\tablerowrule
共享 & 两个程序想读同一份数据，但它们各自记录不同的块号。如何让多个程序引用同一份数据？ & Web 服务器和缓存代理读同一页 \\\tablerowrule
隔离 & 程序 A 不应该能覆盖程序 B 的数据。裸磁盘上任何程序都能写任何块。 & 恶意程序覆盖块 0（引导扇区） \\\tablerowrule
崩溃 & 断电时写操作可能只完成了一半。块 100 写好了，块 101 没写。重启后数据不一致。 & 写到一半断电 \\\tablerowrule
\end{longtable}

前六项说明裸磁盘缺少稳定的对象组织；即使对象已经能够命名和分配，运行期间仍要处理性能、并发与权限问题：

\begin{longtable}{|F{0.14\linewidth}|F{0.58\linewidth}|F{0.22\linewidth}|}
\hline
\rowcolor{BookBlueTint}
\textbf{问题} & \textbf{具体困难} & \textbf{举例} \\\tablerowrule
\endhead
缓存 & 同一个块被反复读取时没有缓存层。每次都走完整 DMA 路径。 & 反复读同一文件的头部 \\\tablerowrule
并发 & 两个程序同时写同一个块，后写的覆盖先写的，数据丢失。 & 两个进程同时 \code{write()} 同一偏移 \\\tablerowrule
碎片化 & 空闲块散布在各处，新文件的数据块不连续，导致随机 I/O 增多。 & 删删写写后空间碎片化 \\\tablerowrule
访问控制 & 谁能读哪些块？裸磁盘没有权限概念。 & 用户 A 读取用户 B 的邮件 \\\tablerowrule
持久化顺序 & 断电后哪些写会"活下来"？DMA 顺序 $\ne$ 发出 \code{write()} 的顺序。 & 先写的数据反而丢了 \\\tablerowrule
\end{longtable}

这 11 个问题并非彼此独立。例如，崩溃恢复需要知道哪些
写已经落盘（持久化顺序），而并发写需要某种锁机制，锁机制又需要
命名来标识被锁的对象。

文件系统（file system）就是处理这些问题的\emph{软件层}。它在裸磁盘
之上构建了一组抽象：文件（命名的字节序列）、目录（命名的层级结构）、
权限（访问控制）、事务或 COW 提交协议（崩溃一致性）以及缓存（性能）。
这里不把“原子写”列成文件系统自动提供的通用属性：不同系统调用、不同文件系统、
不同挂载模式所承诺的可见性与持久性边界并不相同，后文会逐项拆开。

文件系统把这些职责分配给不同结构和协议：超级块提供布局入口，inode 保存文件
属性和块映射，目录项保存名字关系，位图管理空间；分配、查找、缓存和恢复算法再
维持它们之间的一致性。后面的设备路径将说明这些磁盘结构最终怎样转化为 I/O 请求。

%%% ================================================================
%%% Section 2: 块设备接口：从 CPU 到磁盘字节
%%% ================================================================

\fsdivision{读写请求如何到达介质}

\fstopic{块设备作为文件}

Linux 统一了设备接口：一切皆文件。磁盘 /dev/sda 也是一个文件。当你执行：

\begin{lstlisting}[language=C]
int fd = open("/dev/sda", O_RDWR);
\end{lstlisting}

你得到一个文件描述符，它支持 \code{read}、\code{write}、\code{lseek}
和部分同步操作。但它有固定容量，而且 I/O 必须服从设备拓扑和访问模式的约束：

\begin{itemize}\tightlist
\item buffered I/O 可以接受字节范围，内核再通过缓存、合并和拆分生成满足设备限制的请求；\code{O\_DIRECT} 等路径通常要求地址、偏移和长度满足对齐约束，否则可能返回 \code{EINVAL}。
\item 小于物理写入粒度或跨越不利边界的写，可能在块层、设备固件或闪存 FTL 中触发读-修改-写，但这不是每次短写都由 Linux 内核固定执行的步骤。
\item 同步操作要区分页缓存写回、块层 flush/FUA 和设备持久化承诺；具体语义依文件类型、驱动和设备能力而定。
\item 块设备的容量就是可寻址范围。到达容量边界后读取会得到 EOF 或短读，越界写会失败；它与普通文件的区别是容量通常不能通过普通 \code{write} 自动增长。
\end{itemize}

\fstopic{Linux 块 I/O 栈}

从用户空间的 \code{read()} 系统调用到磁盘上的物理扇区，数据要穿过
多层软件栈。理解这个栈对理解文件系统的性能行为至关重要。

\begin{fspdfdiagram}
  \node[os compute, minimum width=11cm, minimum height=0.8cm] (app) at (0,8) {Application: \texttt{read(fd, buf, 4096)}};
  \node[os control, minimum width=11cm, minimum height=0.8cm] (vfs) at (0,7) {VFS: system call entry, fd lookup};
  \node[os compute, minimum width=11cm, minimum height=0.8cm] (fs) at (0,6) {Filesystem (ext4): logical block $\to$ physical block};
  \node[os memory, minimum width=11cm, minimum height=0.8cm] (cache) at (0,5) {Page Cache: page index = offset / 4096};
  \node[os compute, minimum width=11cm, minimum height=0.8cm] (bio) at (0,4) {Block I/O Layer: \texttt{bio\_alloc}, \texttt{submit\_bio}};
  \node[os control, minimum width=11cm, minimum height=0.8cm] (sched) at (0,3) {I/O Scheduler: merge and sort};
  \node[os compute, minimum width=11cm, minimum height=0.8cm] (drv) at (0,2) {Device Driver: DMA descriptor programming};
  \node[os small block, minimum width=11cm, minimum height=0.8cm] (hw) at (0,1) {Hardware Controller: DMA transfer};
  \node[os small block, minimum width=11cm, minimum height=0.8cm] (disk) at (0,0) {Disk: sectors on platter / NAND pages};

  \draw[os bus] (app) -- (vfs);
  \draw[os bus] (vfs) -- (fs);
  \draw[os bus] (fs) -- (cache);
  \draw[os bus] (cache) -- (bio);
  \draw[os bus] (bio) -- (sched);
  \draw[os bus] (sched) -- (drv);
  \draw[os bus] (drv) -- (hw);
  \draw[os strong bus] (hw) -- (disk);

  \node[os label, anchor=west] at (6, 5) {cache miss $\to$};
  \node[os label, anchor=west] at (6, 4) {bio\_alloc};
  \node[os label, anchor=west] at (6, 3) {submit\_bio};
  \node[os label, anchor=west] at (6, 2) {DMA setup};
  \node[os label, anchor=west] at (6, 1) {MSI-X IRQ};
  \node[os label, anchor=west] at (6, 0) {physical media};
\end{fspdfdiagram}

\fsdiagramnote{数据从上到下穿过 8 层。上半部分（VFS、filesystem、page cache）
在内核内存中操作页和 inode；下半部分（bio、scheduler、driver、hardware）
操作扇区和 DMA 描述符。cache miss 是分水岭：命中时数据从内存返回，
miss 时才需要走完整 I/O 路径。}

下面逐层分析。

\fstopic{VFS 层}

VFS（Virtual File System）是所有文件系统的统一入口。当你调用
\code{read(fd, buf, 4096)}：

\begin{enumerate}\tightlist
\item 内核通过 \code{fd} 在当前进程的文件描述符表中找到 \code{struct file}。
\item \code{struct file} 包含一个指向 \code{struct inode} 的指针和当前
  文件偏移量 \code{file->f_pos}。
\item VFS 检查偏移量是否超出文件大小。如果是，返回 0（EOF）。
\item VFS 调用 \code{file->f_op->read} 或 \code{file->f_op->read_iter}，
  这是由具体文件系统（如 ext4）注册的回调。
\end{enumerate}

VFS 本身不碰磁盘——它只是路由。

\fstopic{文件系统层：逻辑块到物理块的映射}

ext4 文件系统收到 VFS 传来的请求后，需要把文件内的\emph{逻辑块号}
映射为磁盘上的\emph{物理块号}。例如，用户想读文件偏移 8192 处的
4096 字节：

\begin{enumerate}\tightlist
\item 逻辑块号 = 8192 / 4096 = 2。
\item ext4 查询 inode 的 \code{i_block[15]} 数组，找到逻辑块 2 对应的
  物理块号（假设为 5000）。
\item 物理扇区号 = 5000 $\times$ 8 = 40000（每个 4K 块 = 8 个 512B 扇区）。
\end{enumerate}

这个映射可能需要多次磁盘读（如果需要遍历间接指针，见
\ref{sec:inode-deep} 节），也可能只需要一次内存查找（如果 inode
已在内存中）。

\fstopic{页缓存层}

Linux 的页缓存（page cache）是文件数据的内存缓存。每个文件的数据
以 4096 字节的页为单位缓存。页缓存的查找键是 \code{(inode, page_index)}。

当文件系统需要读取文件的某页时：

\begin{enumerate}\tightlist
\item 先在页缓存中查找 \code{find_get_page(mapping, index)}。
\item 如果找到且页标记为 \code{PG_uptodate}（数据有效），直接从内存
  复制到用户缓冲区——\emph{完全不碰磁盘}。
\item 如果没找到（cache miss），分配一个新页，标记为 \code{PG_locked}，
  然后向块 I/O 层提交读取请求。
\item 当 I/O 完成后，页被标记为 \code{PG_uptodate} 并解锁。等待的
  \code{read()} 系统调用被唤醒，从页缓存复制数据到用户缓冲区。
\end{enumerate}

Linux 的 \emph{buffered I/O} 以页缓存为默认数据路径：普通 \code{read()} 先查缓存，
miss 时才走设备 I/O；普通 \code{write()} 先修改缓存 folio 并标脏，再由
writeback 机制异步回写。\code{O\_DIRECT} 和 DAX 等路径可以绕过传统页缓存，
所以不能把“buffered I/O 的默认路径”写成“所有文件读写的唯一通路”。
\code{write()} 返回也不等于掉电持久；可靠程序还必须检查 \code{fsync()} 的
返回值，并确认底层设备履行 flush/FUA 等持久化承诺。

\fstopic{块 I/O 层：bio 结构}

当页缓存 miss 时，文件系统需要从磁盘读取数据。它通过构造一个
\code{bio}（block I/O）对象来描述这次 I/O 请求。

\begin{lstlisting}[language=C]
/* simplified block-layer request */
struct bio {
    sector_t           bi_sector;    /* starting sector on device   */
    unsigned int       bi_size;      /* total bytes in this I/O      */
    struct block_device *bi_bdev;    /* target block device          */
    struct bio_vec     *bi_io_vec;   /* array of page fragments      */
    unsigned short     bi_vcnt;      /* count of bio_vecs            */
    unsigned short     bi_idx;       /* current index into bi_io_vec */
    unsigned int       bi_opf;       /* operation flags (READ/WRITE)*/
    bio_end_io_t      *bi_end_io;    /* completion callback          */
    void              *bi_private;   /* caller-private data          */
};

struct bio_vec {
    struct page      *bv_page;      /* page in page cache            */
    unsigned int      bv_len;       /* length of this segment       */
    unsigned int      bv_offset;     /* offset within the page       */
};
\end{lstlisting}

一个 \code{bio} 可以包含多个 \code{bio_vec}，每个 \code{bio_vec} 指向
页缓存中的一个页面片段。这种设计允许一次 I/O 操作将磁盘上连续的扇区
\emph{散布}（scatter）到内存中不连续的多个页面，或反向收集（gather）。

\fssubtopic{一个 1 MiB 读取如何变成多个 bio}

用户调用 \code{read(fd, buf, 1048576)}（读 1 MiB）。假设文件偏移从
0 开始，块大小 4096，页大小 4096：

\begin{enumerate}\tightlist
\item 1 MiB = 1048576 字节 = 256 个页。
\item VFS 为每个页检查页缓存。假设全部 miss。
\item 文件系统为每个页生成一个 \code{bio_vec}：\code{bv_page} 指向
  页缓存中新分配的页，\code{bv_len} = 4096，\code{bv_offset} = 0。
\item ext4 可能将逻辑上连续的多个页合并到一个 \code{bio} 中
  （前提是它们映射到磁盘上连续的物理块）。假设文件在磁盘上连续存放，
  256 个页可以合并成 1 个 \code{bio}：\code{bi_sector} = 起始扇区，
  \code{bi_size} = 1048576，\code{bi_vcnt} = 256。
\item 如果文件在磁盘上不连续（例如跨了两个物理区域），则拆分为
  2 个 \code{bio}。
\end{enumerate}

\fstopic{I/O 调度器：合并与排序}

\code{bio} 被 \code{submit_bio()} 提交到块设备的 I/O 队列后，由
I/O 调度器（I/O scheduler）处理。调度器的核心任务：

\begin{itemize}\tightlist
\item \textbf{合并（merging）}：如果新请求的扇区与队列中已有请求的扇区
  相邻，合并为一个更大的请求。减少请求开销。
\item \textbf{排序（sorting）}：对队列中的请求按扇区号排序，使磁头
  按顺序扫过盘面，减少寻道时间。
\end{itemize}

不同设备需要不同策略：

\begin{longtable}{|F{0.18\linewidth}|F{0.29\linewidth}|F{0.22\linewidth}|F{0.25\linewidth}|}
\hline
\rowcolor{BookBlueTint}
\textbf{调度器} & \textbf{策略} & \textbf{适用设备} & \textbf{关键特性} \\\tablerowrule
\endhead
none & 不调度，直接提交 & NVMe SSD & 随机访问延迟一致，无需排序 \\\tablerowrule
deadline & 读写分离队列，按扇区排序 + 截止时间 & HDD, 通用 & 防止请求饥饿 \\\tablerowrule
cfq（历史） & 单队列时代按进程分配时间片 & 旧式 HDD 桌面系统 & 已退出主线，不作为当前默认策略 \\\tablerowrule
bfq & 预算公平队列，权重分配 & HDD, 桌面/移动 & 交互响应好 \\\tablerowrule
mq-deadline & 多队列版 deadline & NVMe + HDD & 兼顾截止时间与多队列并行 \\\tablerowrule
\end{longtable}

\begin{warningbox}
NVMe 没有机械寻道，因此按扇区排序的收益通常远小于 HDD，但不能据此规定“永远使用 \code{none}”。\code{none}、\code{mq-deadline} 或 BFQ 的选择还取决于延迟隔离、公平性、虚拟化层和工作负载。应在目标设备上比较吞吐与尾延迟，而不是把设备类型直接映射成唯一调度器。
\end{warningbox}

\fstopic{DMA 深入：从描述符到中断}

当 I/O 调度器将请求派发给设备驱动后，驱动程序需要设置 DMA 传输。
现代设备常用描述符环、散布-收集列表或 submission/completion queue。以队列式设备为例：

\begin{enumerate}\tightlist
\item CPU 在内存中构造一个或多个 DMA 描述符。每个描述符包含：
  \begin{itemize}\tightlist
  \item 设备可见的 DMA 地址与长度；启用 IOMMU 时地址通常是 IOVA
  \item 传输长度（字节数）
  \item 标志位（方向、中断使能）
  \item 下一个描述符索引、地址或队列位置，具体由设备格式定义
  \end{itemize}
\item 驱动把描述符发布到设备约定的内存队列；有的设备在初始化时配置队列基址，有的设备使用链式描述符。
\item 驱动执行必要的内存屏障，再更新队列尾指针或写 doorbell，触发控制器处理。
\item 控制器\emph{自行}通过总线读取内存中的描述符（这是 DMA 本身的
  第一次使用——用 DMA 来读 DMA 描述符！）。
\item 控制器按照描述符链执行数据传输：直接在磁盘和 RAM 之间搬移数据，
  \textbf{不经过 CPU 寄存器}。
\item 每个描述符完成后，控制器可以发出一个中断；或者等所有描述符
  完成后发一个中断。现代设备通常用 MSI-X（Message Signaled Interrupts
  with eXtended），设备发起约定地址与数据格式的内存写事务，由平台中断控制器路由到目标 CPU，无需传统中断引线。
\item CPU 收到中断，调用驱动注册的中断处理函数。处理函数标记
  \code{bio} 为完成状态，调用 \code{bi_end_io} 回调。
\end{enumerate}

DMA 允许控制器以总线主控（bus master）身份访问已经映射的内存，不再依赖 CPU
逐字节搬运。CPU 准备描述符并通知设备，控制器执行传输，完成事件再由中断或轮询
路径收取。队列由此可以容纳多个在途请求，CPU 无须为每个数据字节参与传输。

\fstopic{完整追踪：read(fd, buf, 4096) 在偏移 8192 处}

让我们追踪一次完整的 \code{read()} 调用，从用户空间到物理磁盘字节。
假设文件 \code{fd} 对应 inode 中记录逻辑块 2 $\to$ 物理块 5000，
块大小 4096，磁盘扇区 512 字节。

\begin{lstlisting}
Step 1: VFS
  - read(fd, buf, 4096) enters kernel via syscall
  - fd -> file struct, file->f_pos = 8192
  - file size check: 8192 + 4096 <= file_size? yes
  - call file->f_op->read_iter()

Step 2: Page Cache lookup
  - page_index = 8192 / 4096 = 2
  - find_get_page(mapping, 2) -> NULL (cache miss)
  - alloc new page, set PG_locked, add to page cache

Step 3: Filesystem block mapping
  - logical_block = 8192 / 4096 = 2
  - lookup inode i_block[]: logical 2 -> physical block 5000
  - sector = 5000 * 8 = 40000 (8 sectors per 4K block)

Step 4: bio construction
  - bio_alloc(bdev, sector=40000, size=4096, op=READ)
  - bio_vec: bv_page = newly allocated page
             bv_len   = 4096
             bv_offset = 0
  - bi_vcnt = 1

Step 5: submit_bio()
  - bio enters device request queue
  - I/O scheduler checks for merge opportunity:
    is there an adjacent pending request? (maybe merge)
    for NVMe: no sort, just submit to SQ (submission queue)
    for HDD: insert into sorted position by sector

Step 6: Device driver
  - driver pops request from queue
  - builds NVMe command or DMA descriptor:
    opcode = READ
    slba   = 40000 (starting LBA)
    nlb    = 7 (number of logical blocks = 8 sectors - 1, 0-based)
    prp1   = DMA address returned by the DMA mapping API
  - writes command to submission queue tail
  - rings doorbell register (SQyTDBL)

Step 7: DMA transfer
  - NVMe controller fetches command from submission queue (via DMA)
  - controller reads 4096 bytes from NAND at LBA 40000
  - controller writes 4096 bytes through the mapped DMA address
  - CPU is free to run other threads during this transfer
  - transfer latency depends on device, queue depth, cache state and workload

Step 8: Interrupt
  - controller writes completion entry to completion queue
  - controller fires MSI-X interrupt vector N
  - CPU enters driver interrupt handler (top half)
  - driver reads completion queue entry, checks status
  - driver calls bio_end_io(bio) -> marks page PG_uptodate, clears PG_locked
  - wakes up the read() that was sleeping on this page

Step 9: Copy to user
  - VFS resumes: page is now valid (PG_uptodate set)
  - copy_to_user(buf, page_address(page) + 0, 4096)
  - read() returns 4096
\end{lstlisting}

注意 Step 7：设备通过 DMA 映射把数据写入内核页，CPU 不执行设备到内存的逐字节复制。Step 9 的 \code{copy\_to\_user} 仍由 CPU 完成；若使用 direct I/O、零拷贝或用户页 pin，数据路径又会不同。因此“DMA 不占用 CPU 搬运设备数据”不能简写成“整条 I/O 路径 CPU 不碰数据”。

Step 8 的完成上下文取决于驱动和配置：NVMe 可以在中断处理路径中处理完成队列，也可以使用 polling 等模式。NAPI 是网络收包的中断缓和与轮询框架，不能作为通用 NVMe 完成机制的名称。

\begin{warningbox}
\code{read()} 成功返回表示请求的字节已经交给调用者；它不建立新的持久化承诺，也不需要随后调用 \code{fsync()} 来让这次读取“更可靠”。持久化问题属于写路径：\code{write()} 返回、脏页写回、块层完成和设备稳定介质承诺是不同边界。页缓存只说明读取可能从内存命中，不能据此推断数据此前是否由本次 I/O 从设备取回。
\end{warningbox}

%%% ================================================================
%%% Section 3: 最小文件系统：四类对象
%%% ================================================================

\fi

\fsdivision{布局字段如何编码为磁盘字节}

前一节没有预先规定“文件系统必须有哪些结构”，而是从尚未解决的问题逐项得到
块、位图、inode、目录项和超级块。现在才把这些职责落实为具体布局。阅读下面的字段时，
应当始终追问它代表哪一项持久事实；如果两个字段同时声称自己是同一事实的权威来源，
更新时就可能产生无法判定的冲突。

\fstopic{磁盘参数}

为了具体化讨论，我们定义一个\emph{教学文件系统}（Minimal FS, MDFS），
运行在一块 1 GiB 磁盘上：

\begin{itemize}\tightlist
\item 磁盘容量：1 GiB = $2^{30}$ 字节 = 1,073,741,824 字节
\item 块大小：4096 字节（4 KiB）
\item 总块数：$2^{30} / 2^{12} = 2^{18} = 262144$ 块
\item inode 大小：256 字节
\item inode 总数：512
\item inode 表大小：$512 \times 256 / 4096 = 32$ 块
\end{itemize}

磁盘布局如下：

\begin{longtable}{|F{0.12\linewidth}|F{0.15\linewidth}|F{0.60\linewidth}|}
\hline
\rowcolor{BookBlueTint}
\textbf{块号} & \textbf{内容} & \textbf{说明} \\\tablerowrule
\endhead
0 & 引导扇区 & 保留，包含引导加载程序。文件系统不使用。 \\\tablerowrule
1 & 超级块 & 文件系统全局元数据：总块数、空闲块数、inode 表位置等。 \\\tablerowrule
2 & inode 位图 & 每位对应一个 inode，1 = 已分配。512 位 = 64 字节，远小于 1 块。 \\\tablerowrule
3--10 & 块位图 & 每位对应一个数据块，1 = 已分配。262144 位 = 32768 字节 = 8 块。 \\\tablerowrule
11--42 & inode 表 & 32 块 $\times$ 4096 / 256 = 512 个 inode。inode N 在块 $11 + N/16$，偏移 $(N\%16) \times 256$。 \\\tablerowrule
43--262143 & 数据块 & 实际存放文件数据。 \\\tablerowrule
\end{longtable}

这个文件系统需要管理\textbf{四类磁盘上的数据结构}：

\begin{enumerate}\tightlist
\item \textbf{超级块}（superblock）：全局元数据，1 块。
\item \textbf{inode 表}（inode table）：每个文件的元数据，32 块。
\item \textbf{目录项}（directory entries）：目录的内容，存在数据块中。
\item \textbf{空闲空间位图}（block bitmap + inode bitmap）：分配状态，$8 + 1 = 9$ 块。
\end{enumerate}

\fstopic{超级块：全局元数据}

\begin{lstlisting}[language=C]
/* Block 1: superblock (64-byte header, rest is padding) */
struct superblock {
    uint32_t magic;               /* 0x4D534653 = "MDFS" magic number   */
    uint32_t total_blocks;        /* 262144                              */
    uint32_t free_blocks;         /* 262100 (initially, after mkfs)    */
    uint32_t total_inodes;        /* 512                                 */
    uint32_t free_inodes;         /* 511 (inode 1 used by root dir)     */
    uint32_t block_size;          /* 4096                               */
    uint32_t inode_size;          /* 256                                */
    uint32_t root_inode;          /* 1 (root directory's inode number)  */
    uint32_t bitmap_start;        /* 2 (block where inode bitmap starts)*/
    uint32_t inode_table_start;   /* 11                                 */
    uint32_t data_start;          /* 43 (first data block)              */
    uint32_t reserved[5];         /* future use                         */
};  /* total: 16 * 4 = 64 bytes, rest of 4096-byte block is zero */
\end{lstlisting}

当 \code{mkfs} 创建文件系统后，块 1 的前 64 字节（小端序）如下：

\begin{lstlisting}
偏移   字节（hex）
0000:  53 46 53 4D  00 00 04 00  FF FF 03 00  00 02 00 00
       |magic----|  |total_blk-|  |free_blk--|  |total_ino|
0010:  FF 01 00 00  00 10 00 00  00 01 00 00  01 00 00 00
       |free_ino-|  |block_sz-|  |inode_sz-|  |root_ino-|
0020:  02 00 00 00  0B 00 00 00  2B 00 00 00  00 00 00 00
       |bitmap---|  |inode_tbl-|  |data_strt|  |reserved0|
0030:  00 00 00 00  00 00 00 00  00 00 00 00  00 00 00 00
       |reserved1|  |reserved2|  |reserved3|  |reserved4|
\end{lstlisting}

逐字段解读：

\begin{itemize}\tightlist
\item \code{magic} = \code{0x4D534653}：小端序读为字节 \code{53 46 53 4D}，
  即 ASCII "MDFS"（反序）。文件系统挂载时检查此值，不匹配则拒绝挂载。
\item \code{total_blocks} = \code{0x00040000} = 262144。
\item \code{free_blocks} = \code{0x0003FFD5} = 262101。初始时 262144 块
  减去 43 块元数据 = 262101。再减去根目录占用的 1 块（块 43）= 262100。
\item \code{total_inodes} = \code{0x00000200} = 512。
\item \code{free_inodes} = \code{0x000001FF} = 511（inode 1 已分配给根目录）。
\item \code{block_size} = \code{0x00001000} = 4096。
\item \code{inode_size} = \code{0x00000100} = 256。
\item \code{root_inode} = 1，\code{bitmap_start} = 2，
  \code{inode_table_start} = 11（\code{0x0000000B}），\code{data_start} = 43
  （\code{0x0000002B}）。
\end{itemize}

\fstopic{inode：文件的元数据}

inode（index node）是文件系统中每个文件的\emph{元数据锚点}。它存储
除文件名以外的所有文件属性。我们的教学文件系统使用 256 字节的
inode（与 ext4 默认值相同），但核心字段布局沿用 ext2 的经典 128 字节
格式：

\begin{lstlisting}[language=C]
/* 256-byte inode (first 128 bytes shown, matching ext2 layout) */
struct mdfs_inode {
    uint16_t i_mode;          /* file type + permission bits    */
    uint16_t i_uid;           /* owner user ID (low 16 bits)   */
    uint32_t i_size;          /* file size in bytes            */
    uint32_t i_atime;         /* last access time (Unix epoch)  */
    uint32_t i_ctime;         /* inode change time             */
    uint32_t i_mtime;         /* data modification time         */
    uint32_t i_dtime;         /* deletion time (0 if alive)    */
    uint16_t i_gid;           /* group ID (low 16 bits)        */
    uint16_t i_links_count;   /* hard link count               */
    uint32_t i_blocks;        /* sectors used (NOT block count!) */
    uint32_t i_flags;         /* ext4 flags (immutable, append, etc.) */
    uint32_t i_osd1;          /* OS-specific data              */
    uint32_t i_block[15];     /* 12 direct + 1 indirect +      */
                              /* 1 double + 1 triple indirect  */
    uint32_t i_generation;   /* inode version (for NFS)       */
    uint32_t i_file_acl;     /* extended attribute block      */
    uint32_t i_dir_acl;       /* high 32 bits of size (dirs)   */
    uint32_t i_faddr;         /* fragment address (obsolete)   */
    uint8_t  i_osd2[12];     /* OS-specific (uid/gid high,    */
                              /* reserved blocks, etc.)        */
    /* --- ext4 extension: bytes 128--255 --- */
    uint32_t i_ctime_extra;   /* sub-second ctime precision    */
    uint32_t i_mtime_extra;   /* sub-second mtime precision    */
    uint32_t i_atime_extra;   /* sub-second atime precision    */
    uint32_t i_crtime;        /* creation time                 */
    uint32_t i_crtime_extra;  /* sub-second crtime precision   */
    uint32_t i_version_hi;    /* high 32 bits of i_generation  */
    uint32_t i_projid;        /* project ID                    */
    uint8_t  i_checksum[8];  /* inode checksum (crc32c)       */
};
\end{lstlisting}

\fssubtopic{inode 的字节级视图}

假设有一个 10000 字节的普通文件，拥有者 uid=0, gid=0，权限 0644，
创建时间戳为 1700000000（2023-11-14 UTC）。文件数据占用 3 个块
（块 68, 69, 70），因为 $\lceil 10000 / 4096 \rceil = 3$。

该 inode（前 64 字节，小端序）如下：

\begin{lstlisting}
偏移   字节（hex）                字段
0000:  A4 81 00 00  10 27 00 00  00 F1 53 65  00 F1 53 65
       |i_mode---|  |i_size---|  |i_atime-|  |i_ctime-|

0010:  00 F1 53 65  00 00 00 00  00 00 01 00  18 00 00 00
       |i_mtime-|  |i_dtime--|  |uid|gid|  |i_links-| |i_blocks|

0020:  00 00 00 00  00 00 00 00  44 00 00 00  45 00 00 00
       |i_flags-|  |i_osd1---|  |block[0]-|  |block[1]-|

0030:  46 00 00 00  00 00 00 00  00 00 00 00  00 00 00 00
       |block[2]-|  |block[3]-|  |block[4]-|  |block[5]-|
\end{lstlisting}

逐字段解读：

\begin{itemize}\tightlist
\item \code{i_mode} = \code{0x81A4}。高 4 位 \code{0x8} 表示
  \code{S_IFREG}（普通文件）。低 12 位 \code{0x1A4} = 0644（\code{rw-r--r--}）。
  完整值：\code{S_IFREG | 0644} = \code{0100644}（八进制）= \code{0x81A4}。
\item \code{i_size} = \code{0x00002710} = 10000。文件精确大小 10000 字节。
\item \code{i_atime} / \code{i_ctime} / \code{i_mtime} = \code{0x6553F100} =
  1700000000，即 2023-11-14T22:13:20Z。
\item \code{i_dtime} = 0（文件未被删除）。
\item \code{i_uid} = 0（root），\code{i_gid} = 0（root）。
\item \code{i_links_count} = \code{0x0001} = 1（1 个硬链接）。
\item \code{i_blocks} = \code{0x00000018} = 24。
  \textbf{注意：不是 3（块数），而是 24（512 字节扇区数）}。
  3 块 $\times$ 8 扇区/块 = 24 扇区。
\item \code{i_block[0]} = \code{0x00000044} = 68（物理块号）。
  \code{i_block[1]} = 69，\code{i_block[2]} = 70。
  \code{i_block[3..14]} = 0（未使用）。
\end{itemize}

\begin{warningbox}
\code{i_blocks} 采用\textbf{512 字节扇区数}作为单位，并可能计入间接指针块、
扩展属性块等元数据；它不等于当前文件系统块的数量。
一个 1 字节文件在 4K 块系统上 \code{i_blocks} = 8（1 块 $\times$ 8 扇区），
而不是 1。读取该字段时应先确认单位，再换算实际占用量。
\end{warningbox}

\fstopic{目录项：文件名的存放}

文件名\emph{不在 inode 中}——文件名存放在目录文件的数据块里，以
\emph{目录项}（directory entry, dirent）的形式组织。ext2/ext4 的
目录项格式如下：

\begin{lstlisting}[language=C]
/* ext4 directory entry (variable length) */
struct ext4_dirent {
    uint32_t inode;        /* inode number of this entry         */
    uint16_t rec_len;      /* total size of this entry (bytes)   */
    uint8_t  name_len;     /* length of filename (bytes)         */
    uint8_t  file_type;    /* file type (see enum below)         */
    char     name[];       /* filename, NOT null-terminated     */
                          /* padded to 4-byte boundary          */
};

/* file_type values */
#define EXT4_FT_UNKNOWN  0
#define EXT4_FT_REG_FILE 1   /* regular file    */
#define EXT4_FT_DIR      2   /* directory       */
#define EXT4_FT_CHRDEV   3   /* character device */
#define EXT4_FT_BLKDEV   4   /* block device    */
#define EXT4_FT_FIFO     5   /* named pipe      */
#define EXT4_FT_SOCK     6   /* socket           */
#define EXT4_FT_SYMLINK  7   /* symbolic link   */
\end{lstlisting}

\code{rec_len} 的设计是 ext2/ext4 目录项的核心：它允许目录项在块内
\emph{紧凑排列}，同时支持\code{rec_len} 覆盖被删除的条目（删除时不移动
后续条目，只把被删条目的 \code{inode} 置 0 并把它的 \code{rec_len} 加到
前一条目）。

\fssubtopic{目录项的字节级视图}

假设根目录（inode 1）的数据块中有一个条目指向文件 \code{"foo"}
（inode 5），以及标准的 \code{"."} 和 \code{".."} 条目：

\begin{lstlisting}
目录数据块（块 43）内容：

偏移  字节（hex）                             解读
0000: 01 00 00 00 0C 00 01 02  2E 00 00 00   inode=1, rec_len=12, namelen=1, type=DIR, "."
000C: 01 00 00 00 0C 00 02 02  2E 2E 00 00  inode=1, rec_len=12, namelen=2, type=DIR, ".."
0018: 05 00 00 00 0C 00 03 01  66 6F 6F 00  inode=5, rec_len=12, namelen=3, type=REG, "foo"
0024: 00 00 00 00 F4 0F 00 00  ...          inode=0 (deleted), rec_len=4012 (rest of block)
\end{lstlisting}

逐条解读：

\begin{itemize}\tightlist
\item 偏移 0x0000：\code{"."} 条目。\code{inode}=1（根目录自身），
  \code{rec_len}=12（8 字节头 + 1 字节名 + 3 字节填充），
  \code{name_len}=1，\code{file_type}=2（DIR），名字 = \code{"."}.
\item 偏移 0x000C：\code{".."} 条目。\code{inode}=1（父目录也是根目录自身），
  \code{rec_len}=12，\code{name_len}=2，\code{file_type}=2（DIR），
  名字 = \code{".."}.
\item 偏移 0x0018：\code{"foo"} 条目。\code{inode}=5，
  \code{rec_len}=12，\code{name_len}=3，\code{file_type}=1（REG\_FILE），
  名字 = \code{"foo"}（\code{0x66 0x6F 0x6F}），填充 1 字节 \code{0x00}。
\item 偏移 0x0024：剩余空间。\code{inode}=0 标记为空闲，
  \code{rec_len}=4012（从 0x0024 到块末 = 0x1000 的剩余字节数）。
\end{itemize}

\fstopic{空闲空间管理：块位图与 inode 位图}

MDFS 使用\emph{位图}（bitmap）管理空闲资源。每个块对应位图中一位：
1 = 已分配，0 = 空闲。

\begin{itemize}\tightlist
\item \textbf{块位图}：262144 个块需要 262144 位 = 32768 字节 = 8 个块
  （块 3--10）。
\item \textbf{inode 位图}：512 个 inode 需要 512 位 = 64 字节，远少于
  1 个块（块 2）。
\end{itemize}

\fssubtopic{块位图在创建 3 个文件后的状态}

假设 \code{mkfs} 后，创建了根目录（占用块 43），然后创建 3 个文件：
\code{/foo}（100 字节，块 68）、\code{/bar}（200 字节，块 69）、
\code{/baz}（50 字节，块 70）。

已分配的块：0--42（元数据）、43（根目录数据）、68、69、70。

位图前 10 字节如下（bit 0 = 块 0，LSB first）：

\begin{lstlisting}
字节 0 (块 0-7):   FF  = 11111111  块 0-7 全部已分配
字节 1 (块 8-15):  FF  = 11111111  块 8-15 全部已分配
字节 2 (块 16-23): FF  = 11111111  块 16-23 全部已分配
字节 3 (块 24-31): FF  = 11111111  块 24-31 全部已分配
字节 4 (块 32-39): FF  = 11111111  块 32-39 全部已分配
字节 5 (块 40-47): 0F  = 00001111  块 40-43 已分配, 44-47 空闲
字节 6 (块 48-55): 00  = 00000000  全部空闲
字节 7 (块 56-63): 00  = 00000000  全部空闲
字节 8 (块 64-71): 70  = 01110000  块 68-70 已分配, 其余空闲
字节 9 (块 72-79): 00  = 00000000  全部空闲
...
\end{lstlisting}

\begin{warningbox}
ext2/ext4 中位图是\textbf{小端序 bit order}：字节内的 bit 0（LSB）
对应最低编号的块。即字节 \code{0x0F} = \code{0b00001111} 表示块
0--3 已分配，块 4--7 空闲。这与大端序的直觉相反。x86 架构下，
\code{testb} 指令测试的是 LSB，与这种 bit order 一致。
\end{warningbox}

\fstopic{完整追踪：创建 /foo（100 字节）}

现在我们追踪一次完整的 \code{create_file} 操作，从路径解析到磁盘写入。

\begin{lstlisting}
=== create_file("/foo", 100 bytes) ===

Step 1: path_resolve("/foo")
  - Start at root inode (superblock.root_inode = 1)
  - Load inode 1 from inode table
  - Root is a directory (i_mode & S_IFDIR), proceed

Step 2: dir_lookup(root_inode=1, "foo")
  - Read root's data block (block 43)
  - Scan dirents: ".", "..", (deleted entry)
  - "foo" not found -> ENOENT (good, we can create it)

Step 3: alloc_inode()
  - Load inode bitmap (block 2)
  - Scan for first 0 bit: bit 5 is free
  - Set bit 5 -> inode 5 is now allocated
  - Mark inode bitmap block dirty

Step 4: init_inode(5)
  - Write inode 5 at block (11 + 5/16) = block 11, offset (5%16)*256 = 1280
  - Fields:
      i_mode       = S_IFREG | 0644 = 0x81A4
      i_uid        = 0
      i_size       = 0  (will be 100 after write)
      i_atime     = i_ctime = i_mtime = current_time
      i_dtime      = 0
      i_gid        = 0
      i_links_count = 1
      i_blocks     = 0  (no data blocks yet)
      i_block[0..14] = 0 (all zero)
  - Mark inode table block 11 dirty

Step 5: alloc_blocks(1)
  - Load block bitmap (blocks 3-10)
  - Scan for first free bit after data_start (block 43)
  - Block 68 is free (bit 68 in bitmap = 0)
  - Set bit 68 -> block 68 allocated
  - Mark block bitmap block dirty
  - Update superblock: free_blocks--

Step 6: set inode 5 block[0] = 68
  - Write i_block[0] = 68 in inode 5's on-disk record
  - Update i_blocks = 8 (1 block * 8 sectors/block)
  - Update i_size = 100
  - Mark inode table block 11 dirty

Step 7: write data
  - Copy 100 bytes of user data into block 68
  - Zero-fill remaining 4096 - 100 = 3996 bytes of block 68
  - Mark block 68 dirty in page cache

Step 8: dir_add_entry(root_inode=1, "foo", inode=5)
  - Read root's data block (block 43, already in page cache)
  - Find the "deleted" entry at offset 0x0024
    (inode=0, rec_len=4012)
  - Split: new entry takes 12 bytes (8 header + 4 name)
    remaining rec_len = 4012 - 12 = 4000
  - Write dirent:
      inode     = 5
      rec_len   = 12
      name_len  = 3
      file_type = 1 (EXT4_FT_REG_FILE)
      name      = "foo\0" (padded to 4 bytes)
  - Update the following entry's rec_len to 4000
  - Mark root's data block (block 43) dirty
  - Update root inode: i_mtime = current_time, i_size may grow
  - Mark root inode dirty

Step 9: mark dirty
  - Inode 5 dirty (metadata changed)
  - Inode 1 (root) dirty (directory modified)
  - Inode bitmap block dirty (bit 5 set)
  - Block bitmap block dirty (bit 68 set)
  - Superblock dirty (free_blocks decremented)
  - Block 68 dirty (file data written)
  - Block 43 dirty (dirent added)

Step 10: writeback
  - Eventually, the writeback mechanism selects dirty folios
  - Flushes all dirty pages to disk via submit_bio()
  - Order may vary (depends on journal mode)
  - On ext4 with journal=data: metadata changes go through journal first
  - On ext4 with journal=ordered: data is written before metadata
  - fsync() waits for all these writes to complete
\end{lstlisting}

\begin{warningbox}
Step 10 是\textbf{最危险的一步}。如果在 Step 8 写入了目录项但 Step 5
的块位图还没落盘，断电后磁盘状态不一致：inode 5 的数据块 68 被标记为
已分配（因为目录项指向 inode 5），但块位图说 68 空闲，导致块泄漏
（leaked block）。反之，如果块位图先落盘但目录项没落盘，则块 68 被标
记为已分配但没有文件引用它——也是泄漏。ext4 用\emph{日志}（journal）
解决这个顺序问题：将元数据变更先写入日志区，原子提交后再应用到
主区域。
\end{warningbox}

\fstopic{磁盘布局图}

\begin{fspdfdiagram}
  \node[os title] at (0, 10.5) {MDFS 磁盘布局（1 GiB, 262144 块）};

  \node[os small block, minimum width=5cm, minimum height=0.7cm, fill=BookRedTint] (b0) at (0, 9.5) {Block 0: 引导扇区};
  \node[os small block, minimum width=5cm, minimum height=0.7cm, fill=BookBlueTint] (b1) at (0, 8.7) {Block 1: 超级块 (64 B header)};
  \node[os small block, minimum width=5cm, minimum height=0.7cm, fill=BookGreenTint] (b2) at (0, 7.9) {Block 2: inode 位图 (64 B used)};
  \node[os small block, minimum width=5cm, minimum height=0.7cm, fill=BookGreenTint] (b3) at (0, 7.1) {Blocks 3--10: 块位图 (8 块, 32768 B)};
  \node[os small block, minimum width=5cm, minimum height=0.7cm, fill=BookAmberTint] (b4) at (0, 6.3) {Blocks 11--42: inode 表 (32 块, 512 inodes)};
  \node[os small block, minimum width=5cm, minimum height=0.7cm, fill=BookAmberTint] (b5) at (0, 5.5) {\ \ inode 1: root dir};
  \node[os small block, minimum width=5cm, minimum height=0.7cm, fill=BookAmberTint] (b6) at (0, 4.7) {\ \ inode 5: /foo};
  \node[os small block, minimum width=5cm, minimum height=0.7cm, fill=BookMuted] (b7) at (0, 3.9) {Block 43: root dir data};
  \node[os small block, minimum width=5cm, minimum height=0.7cm, fill=BookRedTint] (b8) at (0, 3.1) {Block 68: /foo data (100 B)};
  \node[os small block, minimum width=5cm, minimum height=0.7cm, fill=BookMuted] (b9) at (0, 2.3) {Block 69: /bar data (200 B)};
  \node[os small block, minimum width=5cm, minimum height=0.7cm, fill=BookMuted] (b10) at (0, 1.5) {Block 70: /baz data (50 B)};
  \node[os small block, minimum width=5cm, minimum height=0.7cm] (b11) at (0, 0.7) {Blocks 71--262143: 空闲数据块};

  \node[os label, anchor=west] at (3, 8.7) {\texttt{magic=0x4D534653}};
  \node[os label, anchor=west] at (3, 7.1) {262144 bits = 32768 B};
  \node[os label, anchor=west] at (3, 6.3) {512 $\times$ 256 B};
  \node[os label, anchor=west] at (3, 5.5) {\texttt{i\_block[0]=43}};
  \node[os label, anchor=west] at (3, 4.7) {\texttt{i\_block[0]=68}};
  \node[os label, anchor=west] at (3, 3.1) {dirent: "foo" $\to$ inode 5};
\end{fspdfdiagram}

\fsdiagramnote{红色 = 数据块，蓝色 = 超级块，绿色 = 位图，琥珀色 = inode 表，
灰色 = 根目录数据。注意块 43（根目录数据）和块 68（/foo 数据）都在数据区
（块 43+），它们没有特殊的存储位置——文件系统通过 inode 的 \code{i_block[]}
指针找到它们。}

%%% ================================================================
%%% Section 4: inode 的深层结构
%%% ================================================================

\end{filesystemlayout}
